\documentclass[10pt,a4paper]{amsart}
\usepackage[margin=19mm]{geometry}
\usepackage{amsmath,amssymb,mathtools}
\usepackage[scr=rsfs,cal=euler]{mathalfa}
\usepackage{xfrac}
\usepackage[unicode,hidelinks,bookmarks=false]{hyperref}
\newcommand{\ie}{i.e.\ }
\newcommand{\cf}{cf.\ }
\newcommand{\resp}{respectively\ }
\newcommand{\Subsection}{\subsection}
\providecommand{\nobreakdash}{\nobreak\hbox{-}}
\setlength{\emergencystretch}{2em}
\multlinegap0pt
\usepackage{xcolor}
\usepackage{mathtools}
\usepackage{bm}
\usepackage[shortlabels]{enumitem}
\newtheorem{lemma}{Lemma}
\newcommand{\R}{\mathbb{R}}
\title{Quotient-Categorical Representations for Bellman-Compatible Average-Reward Distributional Reinforcement Learning}
\author{Ege C. Kaya, Aliasghar Pourghani, Vijay Gupta, Abolfazl Hashemi}
\date{Source: arXiv:2605.11289v1 (CC BY 4.0)}
\begin{document}
\maketitle

\noindent\emph{Source and licence.} Quotient-Categorical Representations for Bellman-Compatible Average-Reward Distributional Reinforcement Learning. Version arXiv:2605.11289v1 (CC BY 4.0), distributed by arXiv under the Creative Commons Attribution 4.0 International licence, http://creativecommons.org/licenses/by/4.0/, as recorded in the arXiv licence metadata for this exact version and shown on the abstract page https://arxiv.org/abs/2605.11289v1. Full licence notice: \texttt{licenses/arxiv-2605.11289-CC-BY-4.0.txt}.\par\smallskip

\noindent\textit{Editorial scope.} Counted unit: the article's lemma lem:appC-1 (source0.tex lines 1000--1005). The article's complete original proof of it is delivered with it (source0.tex lines 1006--1043, one \texttt{proof} environment of 38 lines). No mathematics is written, repaired or re-derived: the statement and the proof are the article's own lines, in the article's own notation, and no retained line is substituted or re-flowed.  No further source-local context is needed: every cross-reference of every retained line resolves inside the two delivered spans.  Nothing was omitted from any retained span, so the package carries no omission record.  No retained line contains a citation, so the wrapper carries no bibliography and no bibliography entry.  No retained line contains a figure environment, an image-inclusion command or drawing code, so no image is delivered and the package declares no asset dependency.  Editorial formatting only, in addition: an AMS \texttt{amsart} preamble that restates the article's own declarations and macros used by the retained lines, namely the environments \texttt{lemma} on separate counters, together with the standard packages those lines need.  The retained environments print in this document as Lemma 1 in order of appearance, whereas the article prints the same items with the numbering of its own full document.  The complete native source of the article as distributed in the arXiv e-print for this exact version is bundled unchanged as \texttt{sources/200348/source0.tex}.
\begin{lemma}\label{lem:appC-1}
For every $b, b' \in \R$ and all $p, q \in \Delta_d$,
\begin{equation}
\ell^\Theta_\mathrm C(L_b p, L_{b'}q) \le \ell^\Theta_\mathrm C(p, q) + \Delta^{-1/2}\lvert b - b'\rvert.
\end{equation}
\end{lemma}

\begin{proof}
By the triangle inequality and the non-expansiveness of $L_{b'}$,
\begin{equation}
\ell^\Theta_{\mathrm C}(L_b p, L_{b'}q) \le \ell^\Theta_\mathrm C(L_b p, L_{b'}p) + \ell^\Theta_\mathrm C(L_{b'} p, L_{b'} q) \le \ell^\Theta_\mathrm C(L_b p, L_{b'}p) +\ell^\Theta_\mathrm C(p, q).
\end{equation}
It therefore suffices to bound $\ell^\Theta_\mathrm C(L_b p, L_{b'}p)$. Write $p = \sum_{m=1}^d p_m e_m$, where $e_m$ is the $m$th canonical basis vector. Since scalar categorical projection is linear on mixtures,
\begin{equation}
L_t p = \sum_{m=1}^d p_m L_t e_m\qquad\text{for every } t \in \R.
\end{equation}
Applying triangle inequality to the vector of cumulative masses yields
\begin{equation}\label{eq:sub-into-this}
\ell^\Theta_\mathrm C(L_b p, L_{b'}p) \le \sum_{m=1}^d p_m\ell^\Theta_\mathrm C(L_b e_m, L_{b'}e_m).
\end{equation}
Fix $m$ and define
\begin{equation}
v_m(t) := \bigl(F_{L_t e_m}(\theta_k)\bigr)_{k=1}^{d-1} \in \R^{d-1}.
\end{equation}
Since $\ell^\Theta_\mathrm C(u,v) = \sqrt{\Delta}\bigl\lVert (F_u(\theta_k) - F_v(\theta_k))_{k=1}^{d-1}\bigr\rVert_2$, it is enough to control $t \mapsto v_m(t)$. Let $x = \theta_{m+t}$ be the translated atom location. If $x \le \theta_1$, then $L_t e_m = e_1$ and $v_m(t) = \mathbf 1$. If $x \ ge \theta_d$, then $L_t e_m = e_d$ and $v_m(t) = 0$. If instead $x \in [\theta_k, \theta_{k+1}]$ for some $k \in \{1, \dots, d-1\}$, scalar categorical projection places mass
\begin{equation}
\lambda_k(x) := \frac{\theta_{k+1} -x}{\Delta}
\end{equation}
on $\theta_k$ and mass $1 - \lambda_k(x)$ on $\theta_{k+1}$. Hence
\begin{equation}
v_m(t) = \bigl(\underbrace{0,\dots,0}_{k-1 \text{ entries}},\lambda_k(x),\underbrace{1, \dots, 1}_{d-1-k\text{ entries}}\bigr),
\end{equation}
so on this interval $v_m$ is affine and
\begin{equation}
\frac{\mathrm d}{\mathrm dt} v_m(t) = -\frac{1}{\Delta}e_k.
\end{equation}
Therefore $\bigl\lVert\frac{\mathrm d}{\mathrm dt} v_m(t) \bigr\rVert_2 \le \Delta^{-1}$ wherever the derivative exists. Integrating this bound and multiplying by $\sqrt{\Delta}$ from the definition of $\ell^\Theta_\mathrm C$ gives
\begin{equation}
\ell^\Theta_\mathrm C(L_b e_m, L_{b'}e_m) \le \Delta^{-1/2}\lvert b- b'\rvert.
\end{equation}
Substituting this into \eqref{eq:sub-into-this} yields
\begin{equation}
\ell^\Theta_\mathrm C(L_b p, L_{b'}p) \le \Delta^{-1/2}\lvert b- b'\rvert
\end{equation}
\end{proof}

\end{document}
